\documentclass[11pt]{report}
\usepackage[margin=0.59in]{geometry}
\usepackage{graphicx}
\usepackage{tikz}
\usetikzlibrary{calc}
\usepackage{hyperref}


\newcommand{\mediumsize}{\fontsize{15pt}{16pt}\selectfont} 
\begin{document} 
\begin{titlepage}

    \begin{minipage}[t]{0.95\textwidth}
        \hfill
        \raggedleft
        \textbf{Immanuel Dietrich} \\
        Utrecht, Netherlands \\
        \today\\
        \href{mailto:immanueld49@gmail.com}{immanueld49@gmail.com} \\
        +49 1512 9081675 \\
        \url{http://linkedin.com/in/immanuel-dietrich}
    \end{minipage}

\vspace{0.7em}

\raggedright Guten Tag,\\


\vspace{0.7em}

mein Name ist Immanuel, ich bin 26 Jahre alt und habe vor kurzem mein Masterstudium in den Niederlanden (Utrecht) abgeschlossen. 
 
\vspace{0.7em}


Meine Ausbildung ist ein Bachelor in Umweltingenieurwesen, gefolgt von einem Master in Water Science and Management. Beide Abschlussarbeiten hatten den Schwerpunkt Wasseraufbereitung. 
Damit bringe ich sowohl technisch-ingenieurmäßiges Handwerk als auch ein breiteres Verständnis für Wassermanagement mit und das möchte ich in meiner Karriere geltend machen.

\vspace{0.7em}
 
Konkret kann ich mir mehrere Richtungen vorstellen:
 
\vspace{0.3em}
 
\begin{itemize}
    \item \textbf{Planung und Auslegung von Wasserinfrastruktur} -- etwa bei einem Ingenieurbüro oder einem Wasserversorger, wo technisches Fachwissen zu Wasseraufbereitung und ingenieursspezifische Kenntnisse gefragt ist.
    \item \textbf{Betrieb und Überwachung von Wasserversorgungsnetzen} -- operative Rollen bei Wasserversorgern, bei denen Datenanalyse, Prozessverständnis und Entscheidungsfindung im laufenden Betrieb benötigt werden.
    \item \textbf{Gutachten und Regulierung} -- Positionen bei Behörden oder Beratungsunternehmen, die sich mit der Einhaltung gesetzlicher Grenzwerte (z.\,B. EU-Vorgaben) und Wasserqualität befassen.
\end{itemize}
 
\vspace{0.3em}
 
Ich bin für alle drei Richtungen offen und lasse mich gerne von konkreten Möglichkeiten leiten, statt mich schon jetzt auf eine einzige festzulegen. Da ich noch nicht fließend Niederländisch spreche, aber in den Niederlanden leben möchte, 
interessieren mich sowohl Positionen dort als auch in Deutschland nahe der Niederländischen Grenze. Der Kreis Kleve ist dabei vorallem interessant, da ich dort auch Familie im Eck habe. 
 
\vspace{0.7em}
 
Ich habe meinen Lebenslauf beigefügt. Ich hoffe, dass dies Einblick gibt in meine Interessengebiete.
 

\vspace{0.7em}

Herzlichen Dank für das Lesen meiner Vorstellungen.

\vspace{0.7em}

\raggedright Mit freundlichen Grüßen,\\
\textbf{Immanuel Dietrich}

\end{titlepage}
\end{document}
